% Experiments section, four settings complete.  Numbers from
% results/{harness,codex_harness}_model_{a,b}/*/eval_task.json, 2026-11-13 16:40.
% Fragment: no preamble.  Needs \usepackage{booktabs}.
% Rebuild the live records with:  python run_experiments.py --report
%                                 python run_codex.py --report

\section{Experiments}

\subsection{Setup}

Four settings cross two frameworks with two models over the same 32 tasks
(20 public, 12 self-built):

\begin{table}[h]
\centering
\small
\begin{tabular}{@{}lll@{}}
\toprule
 & \textbf{Model A}: \texttt{deepseek-v4-pro-0813} & \textbf{Model B}: \texttt{qwen3-coder-flash} \\
\midrule
Codex CLI (baseline)   & Experiment 1 & Experiment 2 \\
Self-developed harness & Experiment 3 & Experiment 4 \\
\bottomrule
\end{tabular}
\end{table}

Model A is a reasoning model, model B is not. Both are reached through the
same OpenAI-compatible endpoint (DMXAPI). Every setting receives the same task
prompt, generated by \texttt{benchmark.build\_prompt} from the task text, the
interface declaration and a fixed block of conventions. Every setting is scored
by the same code: the program is compiled by CODESYS~V3.5~SP22 on a Windows
host behind an HTTP service, and then run against the task's test cases on the
local ST interpreter. A program that does not compile scores zero cases.

\paragraph{What each framework is allowed to see.}
The self-developed harness compiles every attempt and, once the code compiles,
runs the first test case of the task and returns the failing scan cycles to the
model. It allows three repairs. The remaining cases are held out and used only
for scoring. Codex runs in a fresh empty workspace holding only
\texttt{TASK.md}, with no access to the compiler or to any test case, and is
judged on the file it leaves behind. The held-out pass rate is the fair
comparison between the two frameworks, because neither of them has seen those
cases.

\paragraph{Codex versions.}
Codex 0.167.2 accepts only the Responses wire API. DMXAPI's Responses endpoint
rejects a tool result sent back with \texttt{role:"tool"} for
\texttt{qwen3-coder-flash}
(\texttt{400 ... role must be one of user,assistant,system,function}), so
Experiment~2 cannot run on it. The same model handles a full tool round trip
correctly over Chat Completions, so the limitation is in the endpoint, not the
model. Experiment~2 therefore runs on Codex 0.44.0 with
\texttt{wire\_api="chat"}; Experiment~1 runs on 0.167.2 with
\texttt{wire\_api="responses"}. The two baseline cells use different Codex
builds, and the comparison between them is confounded by the version. Codex
0.44.0 also reports zero token usage against this endpoint, so Experiment~2 has
no token figures.

\subsection{Results}

\begin{table}[h]
\centering
\small
\begin{tabular}{@{}lrrrr@{}}
\toprule
 & \textbf{Exp.~1} & \textbf{Exp.~2} & \textbf{Exp.~3} & \textbf{Exp.~4} \\
 & Codex, A & Codex, B & Harness, A & Harness, B \\
\midrule
Compiled, first attempt      & 31 / 32 & 11 / 32 & 28 / 32 & 10 / 32 \\
Compiled, final              & 31 / 32 & 11 / 32 & 29 / 32 & 20 / 32 \\
Test cases passed            & 88 / 116 & 6 / 116 & 99 / 116 & 21 / 116 \\
\quad rate                   & 75.9\% & 5.2\% & 85.3\% & 18.1\% \\
Held-out cases passed        & 60 / 84 & 5 / 84 & 70 / 84 & 13 / 84 \\
\quad rate                   & 71.4\% & 6.0\% & 83.3\% & 15.5\% \\
Tasks fully passed           & 24 / 32 & 0 / 32 & 24 / 32 & 3 / 32 \\
Case rate, compiled tasks only & 77.9\% & 16.7\% & 92.5\% & 30.0\% \\
\midrule
Input tokens                 & 7{,}609{,}239 & n/a & 67{,}614 & 163{,}666 \\
Output tokens                & 395{,}037 & n/a & 520{,}864 & 59{,}577 \\
Wall-clock time (s)          & 6{,}084 & 4{,}104 & 8{,}500 & 2{,}182 \\
Median time per task (s)     & 164 & 77 & 247 & 74 \\
\bottomrule
\end{tabular}
\caption{All four settings, 32 tasks each. Codex never sees a test case, so
its ``first'' and ``final'' compile rates are the same measurement.}
\label{tab:main}
\end{table}

\begin{table}[h]
\centering
\small
\begin{tabular}{@{}lrrrr@{}}
\toprule
Tasks fully passed & \textbf{Exp.~1} & \textbf{Exp.~2} & \textbf{Exp.~3} & \textbf{Exp.~4} \\
\midrule
Easy (9)       & 7 & 0 & 7 & 2 \\
Medium (13)    & 10 & 0 & 9 & 0 \\
Hard (10)      & 7 & 0 & 8 & 1 \\
\midrule
Public (20)    & 15 & 0 & 14 & 3 \\
Self-built (12) & 9 & 0 & 10 & 0 \\
\bottomrule
\end{tabular}
\caption{Fully passed tasks by difficulty and by source.}
\end{table}

Per-task rows are in \texttt{LIVE\_RECORD.md} (Experiments 3 and 4) and
\texttt{CODEX\_LIVE\_RECORD.md} (Experiments 1 and 2); the per-case tables in
the course template are in \texttt{TEST\_TABLES.md} and
\texttt{CODEX\_TEST\_TABLES.md}. All four files are regenerated from the result
directories after every task.

\subsection{Framework: harness against Codex, model A}

With model A the two frameworks solve the same number of tasks, 24 of 32, but
not the same tasks: 20 are solved by both, 4 only by Codex, 4 only by the
harness. On the held-out cases the harness is ahead, 83.3\% against 71.4\%, and
on compiled tasks its case rate is 92.5\% against 77.9\%. Codex compiles more
often (31 against 29), the harness passes more cases once it compiles.

The cost is where the two differ most. The harness sends 67{,}614 input tokens
for the whole set, a median of 1{,}019 per task. Codex sends 7{,}609{,}239, a
median of 70{,}538 per task and 2{,}087{,}901 for \texttt{M06} alone. Codex
re-sends the whole conversation, including every file it has read and every
command output, on each tool call; the harness keeps one prompt and appends
only the compiler messages or the failing cycles. Per solved task Codex uses
333{,}512 tokens and the harness 24{,}520, a factor of 13.6. Output tokens go
the other way (395{,}037 against 520{,}864), because the harness asks a
reasoning model for full programs on every repair. Codex is faster per task
(median 164\,s against 247\,s), because it never waits on the compile service.

Of the 8 tasks Codex fails, two are not model errors. \texttt{H04} and
\texttt{M03} compile in CODESYS but use constructs the interpreter does not
support (\texttt{UPPER\_BOUND}, and a user-defined array type indexed on
assignment), so all ten of their cases are \texttt{NOT\_RUN} and score zero. On
an ADS backend against a real runtime they would be scored. The harness passed
both tasks: under the two frameworks the same model wrote different code, and
the interpreter's coverage decided the score.

\subsection{An unstated convention}
\label{sec:convention}

Every test case expects an output to change in the same scan cycle as the
transition that causes it. The reference implementations do this by updating
the state first and assigning the outputs from the new state at the end of the
scan. The conventions block in the prompt states edge triggering, timer
semantics and status encoding, but says nothing about output timing.

An implementation that clears its outputs at the top of the scan and sets
them inside the \texttt{CASE} branch reads the requirement just as well, and
shows every output one cycle late. Codex did exactly this on \texttt{E01}: it
compiled at once and failed all six cases, first at cycle~1 with
\texttt{Pump\_Motor} expected TRUE and actual FALSE. The harness's \texttt{E01}
started the same way, saw the failing feedback case and was corrected, and
passed five of six.

So the convention costs the harness a repair round and costs Codex the task.
Some of the gap on held-out cases between the two frameworks is the value of
test feedback in general; this part of it is the value of being told a
convention the prompt left out. The convention is now written into the
case-format appendix (Section~\ref{app:caseformat}) and should be added to
\texttt{GENERAL\_CONVENTIONS} in the prompt before any further run. We did not
change it during these experiments, because every setting has to see the same
prompt.

\subsection{Model: why Qwen's pass rate is so low}

The two models differ in kind, not only in size. \texttt{deepseek-v4-pro} is a
reasoning model: before it writes any code it produces a long hidden chain of
thought, and the endpoint reports those tokens. \texttt{qwen3-coder-flash} is a
non-reasoning model: it emits the program directly, token by token, with no
planning stage. In the harness DeepSeek's median output is 11{,}114 tokens per
call, of which the program itself is a few hundred; Qwen's is 583, essentially
the program and nothing else. Most of what follows comes from that difference.

Under both frameworks Qwen solves almost nothing: 3 of 32 tasks in the harness
and none in Codex, against 24 for DeepSeek in each. Syntax is only part of it.
Once Qwen's code compiles it still passes only 30.0\% of cases (harness) or
16.7\% (Codex), against 92.5\% and 77.9\% for DeepSeek.

\paragraph{1. The first answer is rarely right.}
The clearest single number: in the harness, DeepSeek's \emph{first} attempt
compiled and passed the feedback case on 25 of 32 tasks. Qwen's did so on 5.
A PLC function block is a state machine in which each output depends on the
current state, the inputs, the previous inputs (for edges) and timer state, all
evaluated in one scan. Getting that right means deciding, before writing, what
holds at cycle~0, which branch a given scan enters, and where in the scan each
output is assigned. A reasoning model does that work in its hidden tokens; a
non-reasoning model has to get it right while writing, in one pass, and on these
tasks it mostly does not.

\paragraph{2. Its errors are the ones planning would prevent.}
Of Qwen's 51 failed cases in the harness, 18 fail at cycle~0, the power-up
state, and 13 more at the first cycle where every mismatched output equals the
previous step's expected value, i.e.\ one scan late. The two sets do not
overlap, so 31 of 51 failures (61\%) are initial-state or ordering errors
rather than wrong process logic. Two examples. In \texttt{C05} the edge detector
is \texttt{startEdge := NOT StartButton AND NOT startEdge;}, which is never TRUE
on the cycle the button is pressed, so the whole sequence starts one cycle late.
In \texttt{C03} \texttt{GateClosed} has no initial value, so at power-up neither
gate output is TRUE and no branch of the state machine applies; the gate never
moves. Both are written correctly in the same task by DeepSeek. Neither needs
domain knowledge Qwen lacks, only a check of what happens in the first two
scans, which is exactly the step a non-reasoning model skips.

\paragraph{3. It cannot use feedback to repair.}
The harness gives every model the same repair loop, and the two models use it
very differently. After a failing test case DeepSeek attempted a repair on 3
tasks and fixed all 3; Qwen attempted 4 and fixed none. After a compile error
DeepSeek recovered on 1 of 4 tasks and Qwen on 10 of 22. Qwen fixes isolated
syntax, but a failing scan cycle requires working back from ``\texttt{Pump\_Motor}
expected TRUE at cycle~1, actual FALSE'' to the assignment that caused it, and
Qwen does not do that. 21 of its 32 harness runs ended because a new attempt was
byte-identical to an earlier one: 11 after a compile error and 10 after test
feedback on code that already compiled. In the second group the model held a
compiling program, was shown the failing cycles, and returned the same source.
DeepSeek reached ``compiled and feedback case passed'' on 29 of 32 tasks under
the same loop.

\paragraph{4. A misleading diagnostic is enough to stop it.}
When a body opens with the Pascal keyword \texttt{BEGIN}, CODESYS does not name
\texttt{BEGIN}. It reports ``\,\texttt{';'} expected instead of
\texttt{'X'}'', where \texttt{X} is the identifier that follows, and adds
``code \texttt{'BEGIN;'} is unreachable, is this intended?''.\footnote{The
compiler runs with a Chinese locale; these are translations. The original
messages are kept verbatim in each task's \texttt{result.json}.} The first line
points at the token after the fault. Qwen hit this in 22 of its 70 harness
compile attempts, spread over 17 tasks, and in 6 of its 12 non-compiling tasks
the \texttt{BEGIN} was still there at the end. In Codex 10 of Qwen's 32 programs
contain \texttt{BEGIN}. DeepSeek is not immune: it wrote \texttt{BEGIN} on 3
harness tasks (\texttt{H02}, \texttt{M07}, \texttt{C12}) and did not recover on
any; those three are exactly the tasks it never compiled. So the diagnostic
is a harness problem too; but it costs DeepSeek three tasks and Qwen most of its
compile failures.

\paragraph{5. It does not drive an agent loop.}
Codex is an agent: the model has to decide which command to run, read the
result and decide again. In Codex Qwen left no source file at all on 7 tasks and
a non-compiling one on 14 more. It ran a median of 11 commands per task and up
to 206, against 4 and 33 for DeepSeek, mostly re-reading \texttt{TASK.md} and
re-creating an empty file. Only three tasks scored any case (\texttt{E03} 3,
\texttt{H01} 2, \texttt{C07} 1). Deciding what to do next is itself a planning
task, and a model with no planning stage does it badly. This cell also runs on
the older Codex build, so part of the effect may be the build; the command
counts show the behaviour regardless.

\paragraph{6. The output cap.}
\texttt{max\_tokens} is 1800 for both models. For Qwen it is a hard cap on the
program (a probe stopped at exactly 1800 with \texttt{finish\_reason=length});
DeepSeek's reasoning tokens are not held to it, and the same probe returned
8{,}747 tokens with \texttt{finish\_reason=stop}. This cut Qwen off once
(\texttt{H05}), so it explains one task, not the gap.

\paragraph{Cost.}
Qwen is cheaper per call and more expensive per solution. In the harness it is
3.3 times faster per task (median 74\,s against 247\,s) and writes 8.7 times
fewer output tokens. Per solved task it costs 74{,}414 tokens and 727\,s,
against 24{,}520 tokens and 354\,s for DeepSeek, because most of its prompt
tokens are spent on repair turns that do not converge (2.84 calls per task
against 1.44). On these tasks a reasoning model's long first answer is cheaper
than a fast model's repeated wrong ones.

\paragraph{Summary.}
Qwen's low pass rate is not mainly a gap in ST knowledge: it writes valid ST
often enough. It fails where a task needs reasoning about time in a scan-cycle
program, i.e.\ initial state, edge detection and output ordering, and it cannot
recover from those errors when shown the failing cycle. A reasoning model does
that work before it writes, which on this benchmark is worth 24 tasks against 3.

\subsection{Measurement issues}

\begin{itemize}
\item \textbf{Codex version.} Experiments 1 and 2 use different Codex builds,
      for the reason given in the setup. Experiment~1 against Experiment~3 is a
      clean framework comparison; Experiment~2 against Experiment~4 is not.
\item \textbf{Output budget.} \texttt{max\_tokens} is 1800 for both models in
      the harness. For Qwen this is a hard cap (a probe stopped at exactly 1800
      with \texttt{finish\_reason=length}); DeepSeek's reasoning tokens are not
      held to it and the same probe returned 8{,}747 tokens with
      \texttt{finish\_reason=stop}. In the recorded runs the cap cut off Qwen
      once, on \texttt{H05}.
\item \textbf{Interpreter coverage.} The ST interpreter does not support
      writing a timer's \texttt{.IN}, \texttt{UPPER\_BOUND}, or indexed assignment
      into a user-defined array type. This affects Qwen's \texttt{E04} in the
      harness and DeepSeek's \texttt{H04} and \texttt{M03} in Codex: 14 cases
      in compiled code that were never executed and score zero.
\item \textbf{Missing \texttt{END\_FUNCTION\_BLOCK}.} The compile service
      treats end of file as the end of the POU, so a program truncated before
      \texttt{END\_FUNCTION\_BLOCK} can still compile. Two of DeepSeek's passing
      harness programs end that way.
\item \textbf{Token figures for Experiment~2} are not available; Codex 0.44.0
      receives zero usage from the endpoint.
\end{itemize}

\subsection{Other work}

\begin{itemize}
\item \textbf{Dataset and specifications.} 12 self-built tasks
      (\texttt{C01}--\texttt{C12}), each with an interface, test cases in the
      unified format and three LTL properties. The 20 public tasks carry 264
      LTL properties (Appendix~\ref{app:public_ltl}). The 22 invariant
      properties of the self-built tasks were checked cycle by cycle against
      the reference implementation and hold; the 14 with \texttt{F} or
      \texttt{X} cannot be decided on finite traces.
\item \textbf{Tooling regression,} as evidence that the measurement works
      rather than as an experimental result: 88/88 public cases on the
      reference ST, 87/88 on the repository ST (one known defect in
      \texttt{H05\_MechArm}, TC07), 40/40 ST mutants and 28/28 Python mutants
      detected, 28/28 self-built cases, and \texttt{evaluate.py --oracle}
      116/116.
\end{itemize}

\subsection{Not done}

\begin{itemize}
\item PLCVerif. The LTL specifications are complete; running them through the
      tool is the remaining step for the verification bonus.
\item A rerun with the output-timing convention in the prompt, which would
      separate the value of the convention from the value of test feedback.
\item Codex with test feedback (\texttt{--with-check}). This would put the two
      frameworks on equal information. It needs the newer compile service
      that supports \texttt{CASE} in the interpreter on the Windows host.
\item Second platform (TIA Portal).
\end{itemize}
